\section{Protocol details}
\label{sec:protocol}

\subsection{Keeping the writer blind to future tasks}
\label{sec:taskblind}
Memory formation must be a function of the past episode alone, $m = W(e)$.
If the writer could see a future task, a positive $\Delta U$ would show only
that a model can be told the answer and copy it forward. We enforce this with
four structural barriers rather than by convention. (1)~\emph{Arity}: the
writer function takes a segment and nothing else. (2)~\emph{A closed field
set}: the writer emits only the fields of its schema. (3)~\emph{A metadata
denylist}: the weakest barrier, present to make a plumbing mistake fail loudly.
(4)~\emph{A single interpolation point}: the prompt has one placeholder, the
segment text. Each barrier carries a negative control that plants the defect
and requires the check to fire. Prompt identity is derived from the prompt
file's bytes, not from a hand-maintained version constant that could report a
stale version.

\subsection{Target position is held fixed}
\label{sec:position}
The \textsc{with} arm holds one more item than the \textsc{without} arm, so
some rule must place it, and that rule is not neutral. Holding the memory,
question, filler, model, decoding and seed fixed and varying only the
position of the target memory changes the answer in 8 of 12 probe cells. Measured at scale in the second-session
run (Section~\ref{sec:sophia}), position changes $\soPosChgPct\%$ of labels. A
comparison that does not control position therefore measures position
together with presence. We hold position fixed across replicates, vary it
across cells, and record it on every row. This is also what makes the
second-reader comparison (\main{sec:readertransfer}) a comparison of readers
and not of positions. The probe is small ($n{=}12$), and we treat the rate
as provisional.

\subsection{The bearing screen and both gold-free screens}
\label{sec:goldfree}
The bearing screen reads the memory text, the task's gold answer, and the
task's evaluation specification; it never reads the question. A pair is
\textsc{bearing} only when the benchmark's scorer marks the memory text
correct against the gold answer \emph{and} the test fires on at most eight of
the 865--1{,}024 units in the task's domain pool. The second clause matters.
An earlier, non-discriminative version of the flag fired on nearly every
candidate: its rate on paired versus unpaired memories was in a ratio of
$0.995$, so it said nothing about any one of them.

We pre-registered two gold-free screens before execution. Each reads the
memory, the question and the same-domain pool, and none of the gold answer,
the reward, any model output, or the evaluation specification. We report
both, because reporting only the screen that worked would select an
instrument on its outcome.
\begin{itemize}
\item \emph{Model-based screen.} The judge is $\lfGfJudge$, the reader's own
model, at temperature 0, asked whether the memory contains the specific
information the question asks about. On cells the gold screen never
selected, its picks move the reader from $\gfWithoutPct\%$ correct without the
memory to $\gfWithPct\%$ with it ($\gfN$ rows, $\gfCells$ cells); its matched
control moves from $\gfCtlWithoutPct\%$ to $\gfCtlWithPct\%$ ($\gfCtlN$ rows,
$\lfGfCtlCells$ cells). The contrast of $\gfDiff$ has a task-clustered interval
of $[\lfGfTaskLo, \lfGfTaskHi]$ ($\lfGfTaskN$ tasks), an episode-clustered one
of $[\lfGfEpLo, \lfGfEpHi]$ ($\lfGfEpN$ episodes), and a paired-by-memory
difference of $\lfGfMemD$ over $\lfGfMemN$ memories. Against the gold screen it
recovers $\lfGfBoth$ of $\lfGfGold$ gold-bearing pairs ($\lfGfRec\%$ recall),
and $\lfGfBoth$ of its $\lfGfPicks$ picks are gold-bearing ($\lfGfPrec\%$). The
run is in one session, separate from the gold run.
\item \emph{Lexical screen.} It returns no effect. A diagnostic independent of
any utility explains why: its bearing cells fail in both arms at a rate of
$\gfLexBaf$, \emph{worse} than unscreened random pairing, against
$\gfLlmBaf$ for the model-based screen. A criterion whose selections are
harder than a random draw has not located bearing pairs.
\end{itemize}

\subsection{Leakage audit}
\label{sec:leakage}
The writer invariant excludes \emph{task-directed writing}, not a naturally
occurring memory that happens to contain a later answer. Such a memory is
genuinely useful, and removing it would condition the cell set on the
outcome. We therefore audit every cell rather than filter it. The detector
reuses the benchmark's own scorer and asks whether the rendered memory text
alone would score correct against the gold answer. For \textsc{bearing} cells
the answer is yes by construction, since that is half the bearing definition.
The audit carries information for the control and \textsc{random} strata and
for the Summary and Structured texts, where containment is not guaranteed and
can be conditioned on in analysis. It is validated by 17
controls (plants it must catch, clean cases it must not fire on, and two
pinned known false negatives) and aborts if any control fails.

\subsection{Three kinds of zero in the cost ledger}
\label{sec:zeros}
Raw makes zero writer calls; Summary and Structured make one each per
segment. The ledger distinguishes three zeros that are routinely merged: a
measured zero (Raw makes zero extraction calls), an absent measurement (no
call was made, so the cell leaves that denominator), and an unreported field
(the endpoint exposes no counter). Failed writes count: a write that failed
after consuming a call still cost that call.
